\section{Docker Registry}

\begin{frame}
  \frametitle{O que é um Registry?}
  Local centralizado para armazenar e compartilhar imagens de container.

  \vspace{0.5em}
  \begin{itemize}
    \item \textbf{Docker Hub}: Registro público padrão, acessível a todos
    \item Possível manter registros \textbf{privados} ou \textbf{locais}
  \end{itemize}

  \vspace{0.5em}
  Outros registros populares:
  \begin{itemize}
    \item Amazon Elastic Container Registry (ECR)
    \item Azure Container Registry (ACR)
    \item Google Container Registry (GCR)
  \end{itemize}
\end{frame}

\begin{frame}
  \frametitle{Registry vs Repository}
  \begin{columns}
    \column{0.5\textwidth}
      \textbf{Registry}
      \begin{itemize}
        \item Local centralizado que armazena imagens
        \item Exemplo: Docker Hub
      \end{itemize}

    \column{0.5\textwidth}
      \textbf{Repository}
      \begin{itemize}
        \item Coleção de imagens relacionadas dentro de um registro
        \item Organiza imagens por projeto
        \item Contém uma ou mais imagens
      \end{itemize}
  \end{columns}
\end{frame}

\begin{frame}
  \frametitle{Fluxo: Build, Tag e Push}
  Clonar o projeto e buildar a imagem:

  \vspace{0.3em}
  \texttt{git clone https://github.com/dockersamples/helloworld-demo-node}

  \vspace{0.3em}
  \texttt{docker build -t SEU\_USERNAME/docker-quickstart .}

  \vspace{0.5em}
  Rodar localmente para testar:

  \vspace{0.3em}
  \texttt{docker run -d -p 8080:8080 SEU\_USERNAME/docker-quickstart}

  \vspace{0.5em}
  Criar uma tag e enviar para o Docker Hub:

  \vspace{0.3em}
  \texttt{docker tag SEU\_USERNAME/docker-quickstart SEU\_USERNAME/docker-quickstart:1.0}

  \vspace{0.3em}
  \texttt{docker push SEU\_USERNAME/docker-quickstart:1.0}
\end{frame}